\DocumentMetadata{pdfstandard=a-2b, lang=ja-JP, tagging=off}

% --- jlreq の不具合の回避策 -------------------------------------------
% \DocumentMetadata（TeX Live 2026）で latex-lab が先に読み込まれると、
% jlreq-complements.sty が定義前の内部命令を呼んでエラーになる（217行目 → 251行目で定義）。
% 先に空の命令として定義しておき、jlreq 自身の定義（\cs_set:Nn）で上書きさせる。
\ExplSyntaxOn%
\cs_if_exist:NF \__jlreq_complements_setup_amsthm_theoremstyle:
  { \cs_new:Npn \__jlreq_complements_setup_amsthm_theoremstyle: { } }
\ExplSyntaxOff%
% ----------------------------------------------------------------------

\documentclass[paper=a4paper,fontsize=11pt]{jlreq}

% 共通設定（src/portfolio.sty）
\usepackage{portfolio}

% バージョン情報（make が output/version.tex を生成。無い場合は既定値）
\IfFileExists{output/version.tex}{%
  \input{output/version.tex}% chktex 27 （lint 時点では未生成のため）
}{%
  \newcommand{\docversion}{dev}%
  \newcommand{\doccommit}{unknown}%
}

\newcommand{\doctitle}{TCP/IPネットワーク基礎と通信性能の考察}
\newcommand{\docauthor}{Koichi Kurita}
\title{\doctitle}
\author{\docauthor}
% PDF のプロパティ（PDF/A の XMP メタデータにも反映される）
\hypersetup{pdftitle={\doctitle}, pdfauthor={\docauthor}}
\date{\today\\[0.5em]{\small \docversion（\doccommit）}}

\begin{document}

\maketitle
\tableofcontents
\newpage

% !TEX root = ../main.tex

\section{はじめに}
\label{sec:intro}

本書ではTCP/IPネットワークの基本と通信性能を解説する。
対象とするプロトコルは、TCP~\cite{rfc9293}およびUDP~\cite{rfc768}である。
\Cref{sec:compare}でTCP\index{TCP}とUDP\index{UDP}の特徴を比較し、
\cref{sec:model}で通信遅延の数式モデルを示す。

% !TEX root = ../main.tex

\section{TCPとUDPの比較}
\label{sec:compare}

\subsection{概要}
\label{sec:overview}

TCPとUDPの主な違いを\cref{tab:tcp-udp}に示す。

\begin{table}[htbp]
\centering
\caption{TCPとUDPの比較}
\label{tab:tcp-udp}
% textlint-disable
\begin{tabular}{lcc}
\toprule
項目 & TCP & UDP \\
\midrule
接続方式 & コネクション型 & コネクションレス型 \\
信頼性 & 高 & 低 \\
遅延 & 高 & 低 \\
\bottomrule
\end{tabular}
% textlint-enable
\end{table}

\subsection{通信の流れ}
\label{sec:flow}

TCPとUDPの通信の流れを\cref{fig:sequence}に示す。
TCPはデータを送る前に、3ウェイハンドシェイクで接続を確立する。
クライアントはSYNを送り、サーバからSYN+ACKを受け取ると、ACKと同時にデータの送信を開始する。
そのため、データがサーバに届くまでに約1.5往復分の時間がかかる。
一方、UDPは接続を確立せずにデータを送るため、片道分の時間で届く。

\begin{figure}[htbp]
\centering
% !TEX root = ../main.tex
% ==========================================
% 図：TCP と UDP の通信シーケンス
% ==========================================
\begin{tikzpicture}[
    font=\small,
    x=0.95cm,
    node/.style={draw, minimum width=16mm, minimum height=6mm, fill=gray!10},
    msg/.style={-{Stealth[length=2mm]}, thick},
    data/.style={msg, blue!60!black},
    lbl/.style={font=\scriptsize, fill=white, inner sep=1pt},
  ]
  % ---------- TCP ----------
  \node[font=\bfseries] at (2, 0.8) {TCP};
  \node[node] (tc) at (0, 0) {クライアント};
  \node[node] (ts) at (4, 0) {サーバ};
  \draw[gray] (tc.south) -- ++(0, -5.6);
  \draw[gray] (ts.south) -- ++(0, -5.6);
  \draw[msg]  (0, -0.8) -- node[lbl, sloped, above] {SYN}           (4, -1.6);
  \draw[msg]  (4, -1.8) -- node[lbl, sloped, above] {SYN+ACK}       (0, -2.6);
  \draw[data] (0, -2.8) -- node[lbl, sloped, above] {ACK+データ}    (4, -3.6);
  \draw[msg]  (4, -3.8) -- node[lbl, sloped, above] {ACK}           (0, -4.6);
  \draw[decorate, decoration={brace, amplitude=4pt, mirror}]
    (-0.25, -0.8) -- node[left=5pt, align=right, font=\scriptsize] {接続確立\\（1往復）} (-0.25, -2.6);
  \draw[decorate, decoration={brace, amplitude=4pt}]
    (4.25, -0.8) -- node[right=5pt, align=left, font=\scriptsize] {データ\\到着まで\\約1.5往復} (4.25, -3.6);

  % ---------- UDP ----------
  \node[font=\bfseries] at (9.5, 0.8) {UDP};
  \node[node] (uc) at (7.5, 0) {クライアント};
  \node[node] (us) at (11.5, 0) {サーバ};
  \draw[gray] (uc.south) -- ++(0, -5.6);
  \draw[gray] (us.south) -- ++(0, -5.6);
  \draw[data] (7.5, -0.8) -- node[lbl, sloped, above] {データ} (11.5, -1.6);
  \draw[decorate, decoration={brace, amplitude=4pt}]
    (11.75, -0.8) -- node[right=5pt, align=left, font=\scriptsize] {データ\\到着まで\\片道のみ} (11.75, -1.6);

  % ---------- 時間軸 ----------
  \draw[-{Stealth}, gray] (-1.9, -0.6) -- (-1.9, -5.4);
  \node[gray, font=\scriptsize, above] at (-1.9, -0.6) {時間};
\end{tikzpicture}

\caption{TCPとUDPの通信シーケンス}
\label{fig:sequence}
\end{figure}

\section{数式モデル}
\label{sec:model}

通信全体の遅延$T$は、\cref{eq:delay}のように各遅延の和で表される。

\begin{equation}
  T = T_{\mathrm{prop}} + T_{\mathrm{trans}} + T_{\mathrm{queue}}
  \label{eq:delay}
\end{equation}

ここで、$T_{\mathrm{prop}}$は伝搬遅延、$T_{\mathrm{trans}}$は送信遅延、
$T_{\mathrm{queue}}$はキューイング遅延である。

\subsection{計算例}
\label{sec:latency}

\Cref{tab:conditions}の条件で\cref{eq:delay}を計算した。
遅延の内訳を\cref{fig:latency}に示す。
伝搬遅延は距離を伝搬速度で割った値、送信遅延はパケットサイズを帯域幅で割った値である。
TCPは接続確立のためにSYNとSYN+ACKの1往復が加わる。

% textlint-disable
\begin{table}[htbp]
\centering
\caption{遅延の計算条件}
\label{tab:conditions}
\begin{tabular}{ll}
\toprule
項目 & 値 \\
\midrule
通信距離 & 1{,}000\,km \\
伝搬速度 & $2 \times 10^{8}$\,m/s（光ファイバ中の目安） \\
帯域幅 & 100\,Mbit/s \\
データパケット & 1{,}500\,バイト \\
制御パケット（SYNなど） & 60\,バイト \\
キューイング遅延 & 1\,ms（仮定） \\
\bottomrule
\end{tabular}
\end{table}
% textlint-enable

計算の結果、UDPの遅延は片道分の6.12\,msである。
TCPは接続確立の12.01\,msが加わり、合計18.13\,msとなる。
なお、これらはモデルによる計算値であり、実測値ではない。
グラフは\texttt{data/latency.csv}の値から描画している。

\begin{figure}[htbp]
\centering
% !TEX root = ../main.tex
% ==========================================
% 図：遅延の内訳（data/latency.csv から描画）
% ==========================================
\pgfplotstableread[col sep=comma]{data/latency.csv}\latencytable%
% 合計列を CSV から計算して追加
\pgfplotstablecreatecol[
  create col/expr={\thisrow{handshake} + \thisrow{prop} + \thisrow{trans} + \thisrow{queue}}
]{total}\latencytable%
\begin{tikzpicture}
  \begin{axis}[
      ybar stacked,
      bar width=14mm,
      width=0.75\linewidth,
      height=60mm,
      ymin=0, ymax=24,
      enlarge x limits=0.5,
      xtick=data,
      xticklabels from table={\latencytable}{protocol},
      ylabel={遅延 [ms]},
      legend style={at={(1.03,0.5)}, anchor=west, font=\small, draw=none},
      legend cell align=left,
      reverse legend,
    ]
    \addplot+[fill=gray!40,      draw=black] table[x expr=\coordindex, y=handshake]{\latencytable};
    \addplot+[fill=blue!35,      draw=black] table[x expr=\coordindex, y=prop]{\latencytable};
    \addplot+[fill=orange!60,    draw=black] table[x expr=\coordindex, y=trans]{\latencytable};
    \addplot+[
        fill=green!35, draw=black,
        nodes near coords={\pgfmathprintnumber[fixed, precision=2]{\pgfplotspointmeta}\,ms},
        point meta=explicit,
        every node near coord/.append style={font=\small, anchor=south},
      ] table[x expr=\coordindex, y=queue, meta=total]{\latencytable};
    \legend{接続確立, 伝搬遅延, 送信遅延, キューイング遅延}
  \end{axis}
\end{tikzpicture}

\caption{遅延の内訳（モデルによる計算値）}
\label{fig:latency}
\end{figure}

% !TEX root = ../main.tex

\section{まとめ}
\label{sec:conclusion}

\Cref{sec:compare}で示したとおり、TCPは信頼性重視、UDPは低遅延重視である。
用途に応じた選択が重要である。


\newpage
\printbibliography[title=参考文献]

\newpage
\printindex

\end{document}
